% ---------------------------------------------------------------------------
% resnet_kv.tex -- hand annotation of key-value pairs in
% "Deep Residual Learning for Image Recognition" (arXiv:1512.03385v1).
%
% Source is the PDF only. Conventions are in ANNOTATION.md.
% Sections are generated: python Resnet/build_annotation.py
% Build:  latexmk -pdf -outdir=Output resnet_kv.tex
% ---------------------------------------------------------------------------
\documentclass[10pt]{article}

\usepackage[letterpaper,margin=2.0cm,top=2.2cm,bottom=2.2cm]{geometry}
\usepackage[T1]{fontenc}
\usepackage{lmodern}
\usepackage{amsmath}
\usepackage{array}
\usepackage{longtable}
\usepackage{booktabs}
\usepackage{xcolor}
\usepackage{titlesec}
\usepackage[hidelinks,breaklinks]{hyperref}
\usepackage{url}

\definecolor{rule}{gray}{0.45}
\definecolor{keycol}{HTML}{1F3864}

\newcolumntype{L}[1]{>{\raggedright\arraybackslash\small}p{#1}}

% \leavevmode is load-bearing: a bare \color at the top of a p{} cell is a
% whatsit contributed in vertical mode, which sets the following paragraph a
% full \baselineskip lower than its neighbours.
\newcommand{\usbreak}{\textunderscore\discretionary{}{}{}}
\newcommand{\K}[1]{\leavevmode{\color{keycol}\ttfamily\let\_\usbreak #1}}
\newcommand{\Sd}[1]{\leavevmode{\color{rule}#1}}

\newcommand{\kvhead}[2]{%
  \caption{#1}\label{#2}\\
  \toprule \textbf{Key} & \textbf{Value} & \textbf{Page} \\ \midrule
  \endfirsthead
  \caption[]{#1\ \emph{(continued)}}\\
  \toprule \textbf{Key} & \textbf{Value} & \textbf{Page} \\ \midrule
  \endhead
  \midrule \multicolumn{3}{r}{\footnotesize\itshape continued on next page}\\
  \endfoot
  \bottomrule
  \endlastfoot}

\newenvironment{kvtable}[2]
  {\begin{longtable}{@{}L{4.3cm} L{10.6cm} L{1.3cm}@{}}\kvhead{#1}{#2}}
  {\end{longtable}}

\setlength{\LTcapwidth}{\textwidth}
\renewcommand{\arraystretch}{1.15}
\titleformat{\section}{\normalfont\large\bfseries}{\thesection}{0.6em}{}
\titlespacing*{\section}{0pt}{1.6em}{0.7em}

\title{\textbf{Key--Value Annotation of}\\[0.25em]
       \emph{Deep Residual Learning for Image Recognition}\\[0.4em]
       {\normalsize He, Zhang, Ren \& Sun --- arXiv:1512.03385v1 [cs.CV], 10 Dec 2015}}
\author{}
\date{}

\begin{document}
\maketitle
\vspace{-3.2em}

\noindent\rule{\textwidth}{0.4pt}
\vspace{0.4em}

\noindent
A hand annotation of the key--value pairs stated in the source paper:
\textbf{547 distinct records} across \textbf{51 keys}. Values are read off the
rendered PDF pages; the pipeline's OCR of the same paper is deliberately not
used as a source, so where the OCR is wrong this annotation still agrees with
the paper.

\medskip
\noindent
Keys are Title Case, values are taken as printed, and the unit lives in the key
rather than the value --- the paper's ``a depth of up to 152 layers'' is
recorded as \K{Depth} / 152. Each distinct \emph{(key, value)} pair is recorded
\textbf{once for the whole paper}, so the count is a count of distinct facts
rather than of sentences. The \textbf{Page} column is annotation metadata, not
part of the pair: it records where the pair was first seen.

\medskip
\noindent
Conventions, and the decisions still open --- chiefly that this flat
\emph{key\,:\,value} form cannot express which table row a number belongs
to --- are written up in \texttt{ANNOTATION.md}.

\medskip
\noindent\rule{\textwidth}{0.4pt}

\tableofcontents

\section{Bibliographic}

\begin{kvtable}{Title}{tab:sec_bibliographic-title}
\K{Title} & Deep Residual Learning for Image Recognition & \Sd{p.1}\\
\end{kvtable}

\begin{kvtable}{Author}{tab:sec_bibliographic-author}
\K{Author} & Kaiming He & \Sd{p.1}\\
\K{Author} & Xiangyu Zhang & \Sd{p.1}\\
\K{Author} & Shaoqing Ren & \Sd{p.1}\\
\K{Author} & Jian Sun & \Sd{p.1}\\
\end{kvtable}

\begin{kvtable}{Organization}{tab:sec_bibliographic-organization}
\K{Organization} & Microsoft Research & \Sd{p.1}\\
\end{kvtable}

\begin{kvtable}{Email}{tab:sec_bibliographic-email}
\K{Email} & \{kahe, v-xiangz, v-shren, jiansun\}@microsoft.com & \Sd{p.1}\\
\end{kvtable}

\begin{kvtable}{Identifier}{tab:sec_bibliographic-identifier}
\K{Identifier} & arXiv:1512.03385v1 & \Sd{p.1}\\
\end{kvtable}

\begin{kvtable}{Category}{tab:sec_bibliographic-category}
\K{Category} & cs.CV & \Sd{p.1}\\
\end{kvtable}

\begin{kvtable}{Date}{tab:sec_bibliographic-date}
\K{Date} & 10 Dec 2015 & \Sd{p.1}\\
\K{Date} & 2015-11-26 & \Sd{p.11}\\
\end{kvtable}

\begin{kvtable}{URL}{tab:sec_bibliographic-url}
\K{URL} & \url{http://image-net.org/challenges/LSVRC/2015/} & \Sd{p.1}\\
\K{URL} & \url{http://mscoco.org/dataset/#detections-challenge2015} & \Sd{p.1}\\
\K{URL} & \url{http://host.robots.ox.ac.uk:8080/leaderboard/displaylb.php?challengeid=11&compid=4} & \Sd{p.11}\\
\K{URL} & \url{http://host.robots.ox.ac.uk:8080/anonymous/3OJ4OJ.html} & \Sd{p.11}\\
\end{kvtable}

\begin{kvtable}{Software}{tab:sec_bibliographic-software}
\K{Software} & Caffe & \Sd{p.2}\\
\end{kvtable}

\section{Datasets and tasks}

\begin{kvtable}{Dataset}{tab:sec_datasets-dataset}
\K{Dataset} & ImageNet & \Sd{p.1}\\
\K{Dataset} & CIFAR-10 & \Sd{p.1}\\
\K{Dataset} & COCO & \Sd{p.1}\\
\K{Dataset} & ImageNet 2012 & \Sd{p.4}\\
\K{Dataset} & MNIST & \Sd{p.7}\\
\K{Dataset} & PASCAL VOC 2007 & \Sd{p.8}\\
\K{Dataset} & PASCAL VOC 2012 & \Sd{p.8}\\
\K{Dataset} & MS COCO & \Sd{p.8}\\
\K{Dataset} & ImageNet DET & \Sd{p.11}\\
\end{kvtable}

\begin{kvtable}{Task}{tab:sec_datasets-task}
\K{Task} & Classification & \Sd{p.1}\\
\K{Task} & Object detection & \Sd{p.1}\\
\K{Task} & ImageNet detection & \Sd{p.1}\\
\K{Task} & ImageNet localization & \Sd{p.1}\\
\K{Task} & COCO detection & \Sd{p.1}\\
\K{Task} & COCO segmentation & \Sd{p.1}\\
\K{Task} & image classification & \Sd{p.1}\\
\K{Task} & visual recognition & \Sd{p.1}\\
\K{Task} & image retrieval & \Sd{p.2}\\
\K{Task} & vector quantization & \Sd{p.2}\\
\K{Task} & localization & \Sd{p.12}\\
\end{kvtable}

\begin{kvtable}{Split}{tab:sec_datasets-split}
\K{Split} & 45k/5k train/val split & \Sd{p.7}\\
\K{Split} & VOC 07 test & \Sd{p.8}\\
\K{Split} & VOC 12 test & \Sd{p.8}\\
\K{Split} & 07+12 & \Sd{p.10}\\
\K{Split} & 07++12 & \Sd{p.10}\\
\K{Split} & COCO train & \Sd{p.11}\\
\K{Split} & COCO val & \Sd{p.11}\\
\K{Split} & COCO trainval & \Sd{p.11}\\
\K{Split} & COCO test-dev & \Sd{p.11}\\
\K{Split} & COCO+07+12 & \Sd{p.11}\\
\K{Split} & COCO+07++12 & \Sd{p.11}\\
\K{Split} & val1/val2 & \Sd{p.11}\\
\end{kvtable}

\begin{kvtable}{Size}{tab:sec_datasets-size}
\K{Size} & 1.28 million training images & \Sd{p.4}\\
\K{Size} & 50k validation images & \Sd{p.4}\\
\K{Size} & 100k test images & \Sd{p.4}\\
\K{Size} & 50k training images & \Sd{p.7}\\
\K{Size} & 10k test images & \Sd{p.7}\\
\K{Size} & 5k trainval & \Sd{p.10}\\
\K{Size} & 16k trainval & \Sd{p.10}\\
\K{Size} & 10k trainval+test & \Sd{p.10}\\
\K{Size} & 80k images & \Sd{p.10}\\
\K{Size} & 40k images & \Sd{p.10}\\
\K{Size} & 80k+40k trainval & \Sd{p.11}\\
\K{Size} & 20k test-dev & \Sd{p.11}\\
\end{kvtable}

\begin{kvtable}{Classes}{tab:sec_datasets-classes}
\K{Classes} & 1000 classes & \Sd{p.4}\\
\K{Classes} & 10 classes & \Sd{p.7}\\
\K{Classes} & 80 object categories & \Sd{p.10}\\
\K{Classes} & 200 object categories & \Sd{p.11}\\
\K{Classes} & 1000-class & \Sd{p.12}\\
\end{kvtable}

\section{Architecture}

\begin{kvtable}{Architecture}{tab:sec_architecture-architecture}
\K{Architecture} & deep convolutional neural networks & \Sd{p.1}\\
\K{Architecture} & VGG nets & \Sd{p.1}\\
\K{Architecture} & residual nets & \Sd{p.1}\\
\K{Architecture} & plain & \Sd{p.1}\\
\K{Architecture} & VLAD & \Sd{p.2}\\
\K{Architecture} & Fisher Vector & \Sd{p.2}\\
\K{Architecture} & Multigrid & \Sd{p.2}\\
\K{Architecture} & MLPs & \Sd{p.2}\\
\K{Architecture} & inception & \Sd{p.2}\\
\K{Architecture} & highway networks & \Sd{p.2}\\
\K{Architecture} & ResNet & \Sd{p.2}\\
\K{Architecture} & VGG-19 & \Sd{p.3}\\
\K{Architecture} & 34-layer plain & \Sd{p.4}\\
\K{Architecture} & 34-layer residual & \Sd{p.4}\\
\K{Architecture} & plain-18 & \Sd{p.5}\\
\K{Architecture} & plain-34 & \Sd{p.5}\\
\K{Architecture} & ResNet-18 & \Sd{p.5}\\
\K{Architecture} & ResNet-34 & \Sd{p.5}\\
\K{Architecture} & VGG-16 & \Sd{p.6}\\
\K{Architecture} & GoogLeNet & \Sd{p.6}\\
\K{Architecture} & PReLU-net & \Sd{p.6}\\
\K{Architecture} & ResNet-34 A & \Sd{p.6}\\
\K{Architecture} & ResNet-34 B & \Sd{p.6}\\
\K{Architecture} & ResNet-34 C & \Sd{p.6}\\
\K{Architecture} & ResNet-50 & \Sd{p.6}\\
\K{Architecture} & ResNet-101 & \Sd{p.6}\\
\K{Architecture} & ResNet-152 & \Sd{p.6}\\
\K{Architecture} & BN-inception & \Sd{p.6}\\
\K{Architecture} & ResNet-50/101/152 & \Sd{p.6}\\
\K{Architecture} & bottleneck & \Sd{p.6}\\
\K{Architecture} & VGG-16/19 & \Sd{p.7}\\
\K{Architecture} & Maxout & \Sd{p.7}\\
\K{Architecture} & NIN & \Sd{p.7}\\
\K{Architecture} & DSN & \Sd{p.7}\\
\K{Architecture} & FitNet & \Sd{p.7}\\
\K{Architecture} & Highway & \Sd{p.7}\\
\K{Architecture} & ResNet-20 & \Sd{p.7}\\
\K{Architecture} & ResNet-56 & \Sd{p.7}\\
\K{Architecture} & ResNet-110 & \Sd{p.7}\\
\K{Architecture} & 1202-layer network & \Sd{p.7}\\
\K{Architecture} & Faster R-CNN & \Sd{p.8}\\
\K{Architecture} & Fast R-CNN & \Sd{p.10}\\
\K{Architecture} & ResNet-50/101 & \Sd{p.10}\\
\K{Architecture} & OverFeat & \Sd{p.12}\\
\K{Architecture} & R-CNN & \Sd{p.12}\\
\K{Architecture} & RPN & \Sd{p.12}\\
\end{kvtable}

\begin{kvtable}{Component}{tab:sec_architecture-component}
\K{Component} & shortcut connections & \Sd{p.2}\\
\K{Component} & identity shortcuts & \Sd{p.2}\\
\K{Component} & weight layer & \Sd{p.2}\\
\K{Component} & ReLU & \Sd{p.3}\\
\K{Component} & global average pooling & \Sd{p.3}\\
\K{Component} & 1000-way fully-connected layer & \Sd{p.3}\\
\K{Component} & softmax & \Sd{p.3}\\
\K{Component} & batch normalization (BN) & \Sd{p.4}\\
\K{Component} & projection shortcut & \Sd{p.4}\\
\K{Component} & zero-padding shortcuts & \Sd{p.4}\\
\K{Component} & max pool & \Sd{p.5}\\
\K{Component} & 10-way fully-connected layer & \Sd{p.7}\\
\K{Component} & region proposal network (RPN) & \Sd{p.10}\\
\K{Component} & RoI pooling & \Sd{p.10}\\
\K{Component} & Spatial Pyramid Pooling & \Sd{p.10}\\
\K{Component} & anchor & \Sd{p.12}\\
\K{Component} & cls & \Sd{p.12}\\
\K{Component} & reg & \Sd{p.12}\\
\end{kvtable}

\begin{kvtable}{Depth}{tab:sec_architecture-depth}
\K{Depth} & 152 & \Sd{p.1}\\
\K{Depth} & 100 & \Sd{p.1}\\
\K{Depth} & 1000 & \Sd{p.1}\\
\K{Depth} & 20 & \Sd{p.1}\\
\K{Depth} & 56 & \Sd{p.1}\\
\K{Depth} & sixteen & \Sd{p.1}\\
\K{Depth} & thirty & \Sd{p.1}\\
\K{Depth} & 34 & \Sd{p.3}\\
\K{Depth} & 18 & \Sd{p.5}\\
\K{Depth} & 50 & \Sd{p.6}\\
\K{Depth} & 101 & \Sd{p.6}\\
\K{Depth} & 19 & \Sd{p.7}\\
\K{Depth} & 32 & \Sd{p.7}\\
\K{Depth} & 44 & \Sd{p.7}\\
\K{Depth} & 110 & \Sd{p.7}\\
\K{Depth} & 1202 & \Sd{p.7}\\
\end{kvtable}

\begin{kvtable}{Layers}{tab:sec_architecture-layers}
\K{Layers} & 34 & \Sd{p.3}\\
\K{Layers} & two or three layers & \Sd{p.3}\\
\K{Layers} & 3 layers & \Sd{p.6}\\
\K{Layers} & 6n+2 & \Sd{p.7}\\
\K{Layers} & 3n & \Sd{p.7}\\
\K{Layers} & 1+2n & \Sd{p.7}\\
\K{Layers} & 2n & \Sd{p.7}\\
\K{Layers} & 91 conv layers & \Sd{p.10}\\
\K{Layers} & 13 conv layers & \Sd{p.10}\\
\end{kvtable}

\begin{kvtable}{Filters}{tab:sec_architecture-filters}
\K{Filters} & 3x3 & \Sd{p.3}\\
\K{Filters} & 1x1 & \Sd{p.4}\\
\K{Filters} & 7x7 & \Sd{p.5}\\
\K{Filters} & 64 & \Sd{p.5}\\
\K{Filters} & 128 & \Sd{p.5}\\
\K{Filters} & 256 & \Sd{p.5}\\
\K{Filters} & 512 & \Sd{p.5}\\
\K{Filters} & 1024 & \Sd{p.5}\\
\K{Filters} & 2048 & \Sd{p.5}\\
\K{Filters} & \{16, 32, 64\} & \Sd{p.7}\\
\end{kvtable}

\begin{kvtable}{Output size}{tab:sec_architecture-output-size}
\K{Output size} & 112x112 & \Sd{p.5}\\
\K{Output size} & 56x56 & \Sd{p.5}\\
\K{Output size} & 28x28 & \Sd{p.5}\\
\K{Output size} & 14x14 & \Sd{p.5}\\
\K{Output size} & 7x7 & \Sd{p.5}\\
\K{Output size} & 1x1 & \Sd{p.5}\\
\K{Output size} & 32x32 & \Sd{p.7}\\
\K{Output size} & 16x16 & \Sd{p.7}\\
\K{Output size} & 8x8 & \Sd{p.7}\\
\end{kvtable}

\begin{kvtable}{FLOPs}{tab:sec_architecture-flops}
\K{FLOPs} & 3.6 billion & \Sd{p.3}\\
\K{FLOPs} & 19.6 billion & \Sd{p.3}\\
\K{FLOPs} & 1.8x10\textasciicircum{}9 & \Sd{p.5}\\
\K{FLOPs} & 3.6x10\textasciicircum{}9 & \Sd{p.5}\\
\K{FLOPs} & 3.8x10\textasciicircum{}9 & \Sd{p.5}\\
\K{FLOPs} & 7.6x10\textasciicircum{}9 & \Sd{p.5}\\
\K{FLOPs} & 11.3x10\textasciicircum{}9 & \Sd{p.5}\\
\K{FLOPs} & 3.8 billion & \Sd{p.7}\\
\K{FLOPs} & 11.3 billion & \Sd{p.7}\\
\K{FLOPs} & 15.3/19.6 billion & \Sd{p.7}\\
\end{kvtable}

\begin{kvtable}{Parameters}{tab:sec_architecture-parameters}
\K{Parameters} & 2.5M & \Sd{p.7}\\
\K{Parameters} & 2.3M & \Sd{p.7}\\
\K{Parameters} & 1.25M & \Sd{p.7}\\
\K{Parameters} & 0.27M & \Sd{p.7}\\
\K{Parameters} & 0.46M & \Sd{p.7}\\
\K{Parameters} & 0.66M & \Sd{p.7}\\
\K{Parameters} & 0.85M & \Sd{p.7}\\
\K{Parameters} & 1.7M & \Sd{p.7}\\
\K{Parameters} & 19.4M & \Sd{p.7}\\
\end{kvtable}

\begin{kvtable}{Stride}{tab:sec_architecture-stride}
\K{Stride} & 2 & \Sd{p.3}\\
\K{Stride} & 16 pixels & \Sd{p.10}\\
\end{kvtable}

\section{Method}

\begin{kvtable}{Method}{tab:sec_method-method}
\K{Method} & residual learning framework & \Sd{p.1}\\
\K{Method} & normalized initialization & \Sd{p.1}\\
\K{Method} & intermediate normalization layers & \Sd{p.1}\\
\K{Method} & stochastic gradient descent (SGD) & \Sd{p.1}\\
\K{Method} & backpropagation & \Sd{p.1}\\
\K{Method} & deep residual learning & \Sd{p.2}\\
\K{Method} & hierarchical basis preconditioning & \Sd{p.2}\\
\K{Method} & 10-crop testing & \Sd{p.4}\\
\K{Method} & data augmentation & \Sd{p.7}\\
\K{Method} & maxout & \Sd{p.8}\\
\K{Method} & dropout & \Sd{p.8}\\
\K{Method} & box refinement & \Sd{p.10}\\
\K{Method} & global context & \Sd{p.10}\\
\K{Method} & multi-scale testing & \Sd{p.10}\\
\K{Method} & non-maximum suppression (NMS) & \Sd{p.10}\\
\K{Method} & box voting & \Sd{p.10}\\
\K{Method} & ensemble & \Sd{p.11}\\
\K{Method} & per-class regression (PCR) & \Sd{p.12}\\
\K{Method} & oracle & \Sd{p.12}\\
\K{Method} & 1-crop & \Sd{p.12}\\
\K{Method} & dense & \Sd{p.12}\\
\end{kvtable}

\begin{kvtable}{Problem}{tab:sec_method-problem}
\K{Problem} & vanishing/exploding gradients & \Sd{p.1}\\
\K{Problem} & degradation & \Sd{p.1}\\
\K{Problem} & optimization difficulty & \Sd{p.7}\\
\K{Problem} & overfitting & \Sd{p.8}\\
\end{kvtable}

\begin{kvtable}{Concept}{tab:sec_method-concept}
\K{Concept} & depth & \Sd{p.1}\\
\K{Concept} & low/mid/high-level features & \Sd{p.1}\\
\K{Concept} & identity mapping & \Sd{p.1}\\
\K{Concept} & residual function & \Sd{p.3}\\
\K{Concept} & residual mapping & \Sd{p.3}\\
\K{Concept} & underlying mapping & \Sd{p.3}\\
\K{Concept} & option A & \Sd{p.6}\\
\K{Concept} & option B & \Sd{p.6}\\
\K{Concept} & option C & \Sd{p.6}\\
\K{Concept} & bottleneck design & \Sd{p.6}\\
\end{kvtable}

\begin{kvtable}{Equation}{tab:sec_method-equation}
\K{Equation} & (1) & \Sd{p.3}\\
\K{Equation} & (2) & \Sd{p.3}\\
\end{kvtable}

\begin{kvtable}{Figure}{tab:sec_method-figure}
\K{Figure} & Fig. 1 & \Sd{p.1}\\
\K{Figure} & Fig. 4 & \Sd{p.1}\\
\K{Figure} & Fig. 2 & \Sd{p.2}\\
\K{Figure} & Fig. 3 & \Sd{p.3}\\
\K{Figure} & Fig. 5 & \Sd{p.3}\\
\K{Figure} & Fig. 7 & \Sd{p.3}\\
\K{Figure} & Fig. 6 & \Sd{p.7}\\
\end{kvtable}

\begin{kvtable}{Table}{tab:sec_method-table}
\K{Table} & Table 1 & \Sd{p.4}\\
\K{Table} & Table 2 & \Sd{p.4}\\
\K{Table} & Table 3 & \Sd{p.6}\\
\K{Table} & Table 4 & \Sd{p.6}\\
\K{Table} & Table 5 & \Sd{p.6}\\
\K{Table} & Table 6 & \Sd{p.7}\\
\K{Table} & Table 7 & \Sd{p.8}\\
\K{Table} & Table 8 & \Sd{p.8}\\
\K{Table} & Table 9 & \Sd{p.10}\\
\K{Table} & Table 10 & \Sd{p.11}\\
\K{Table} & Table 11 & \Sd{p.11}\\
\K{Table} & Table 12 & \Sd{p.11}\\
\K{Table} & Table 13 & \Sd{p.12}\\
\K{Table} & Table 14 & \Sd{p.12}\\
\end{kvtable}

\section{Training setup}

\begin{kvtable}{Learning rate}{tab:sec_training-learning-rate}
\K{Learning rate} & 0.1 & \Sd{p.4}\\
\K{Learning rate} & 0.01 & \Sd{p.7}\\
\K{Learning rate} & 0.001 & \Sd{p.10}\\
\K{Learning rate} & 0.0001 & \Sd{p.10}\\
\end{kvtable}

\begin{kvtable}{Weight decay}{tab:sec_training-weight-decay}
\K{Weight decay} & 0.0001 & \Sd{p.4}\\
\end{kvtable}

\begin{kvtable}{Momentum}{tab:sec_training-momentum}
\K{Momentum} & 0.9 & \Sd{p.4}\\
\end{kvtable}

\begin{kvtable}{Batch size}{tab:sec_training-batch-size}
\K{Batch size} & 256 & \Sd{p.4}\\
\K{Batch size} & 128 & \Sd{p.7}\\
\K{Batch size} & 8 images & \Sd{p.10}\\
\K{Batch size} & 16 images & \Sd{p.10}\\
\end{kvtable}

\begin{kvtable}{Iterations}{tab:sec_training-iterations}
\K{Iterations} & 60 x 10\textasciicircum{}4 & \Sd{p.4}\\
\K{Iterations} & 32k & \Sd{p.7}\\
\K{Iterations} & 48k & \Sd{p.7}\\
\K{Iterations} & 64k & \Sd{p.7}\\
\K{Iterations} & 400 iterations & \Sd{p.7}\\
\K{Iterations} & 240k & \Sd{p.10}\\
\K{Iterations} & 80k & \Sd{p.10}\\
\end{kvtable}

\begin{kvtable}{Crop}{tab:sec_training-crop}
\K{Crop} & 224x224 & \Sd{p.4}\\
\K{Crop} & 32x32 & \Sd{p.7}\\
\end{kvtable}

\begin{kvtable}{Scale}{tab:sec_training-scale}
\K{Scale} & [256, 480] & \Sd{p.4}\\
\K{Scale} & \{224, 256, 384, 480, 640\} & \Sd{p.4}\\
\K{Scale} & 600 pixels & \Sd{p.10}\\
\K{Scale} & \{200, 400, 600, 800, 1000\} & \Sd{p.10}\\
\end{kvtable}

\begin{kvtable}{Augmentation}{tab:sec_training-augmentation}
\K{Augmentation} & horizontal flip & \Sd{p.4}\\
\K{Augmentation} & per-pixel mean subtracted & \Sd{p.4}\\
\K{Augmentation} & color augmentation & \Sd{p.4}\\
\K{Augmentation} & 4 pixels are padded & \Sd{p.7}\\
\end{kvtable}

\begin{kvtable}{Hardware}{tab:sec_training-hardware}
\K{Hardware} & two GPUs & \Sd{p.7}\\
\K{Hardware} & 8-GPU implementation & \Sd{p.10}\\
\K{Hardware} & 1 per GPU & \Sd{p.10}\\
\end{kvtable}

\begin{kvtable}{Threshold}{tab:sec_training-threshold}
\K{Threshold} & 0.3 & \Sd{p.10}\\
\end{kvtable}

\begin{kvtable}{Ratio}{tab:sec_training-ratio}
\K{Ratio} & 1:1 & \Sd{p.12}\\
\end{kvtable}

\begin{kvtable}{Proposals}{tab:sec_training-proposals}
\K{Proposals} & 300 proposals & \Sd{p.10}\\
\K{Proposals} & 200 proposals & \Sd{p.12}\\
\end{kvtable}

\begin{kvtable}{Anchors}{tab:sec_training-anchors}
\K{Anchors} & 8 anchors & \Sd{p.12}\\
\end{kvtable}

\section{Results}

\begin{kvtable}{Metric}{tab:sec_results-metric}
\K{Metric} & top-1 error & \Sd{p.4}\\
\K{Metric} & top-5 error & \Sd{p.4}\\
\K{Metric} & training error & \Sd{p.5}\\
\K{Metric} & test error & \Sd{p.5}\\
\K{Metric} & validation error & \Sd{p.5}\\
\K{Metric} & mAP@.5 & \Sd{p.8}\\
\K{Metric} & mAP@[.5, .95] & \Sd{p.8}\\
\K{Metric} & top-5 localization err & \Sd{p.12}\\
\K{Metric} & LOC error on GT CLS & \Sd{p.12}\\
\end{kvtable}

\begin{kvtable}{Error}{tab:sec_results-error}
\K{Error} & 3.57\% & \Sd{p.1}\\
\K{Error} & 27.94 & \Sd{p.5}\\
\K{Error} & 27.88 & \Sd{p.5}\\
\K{Error} & 28.54 & \Sd{p.5}\\
\K{Error} & 25.03 & \Sd{p.5}\\
\K{Error} & 28.07 & \Sd{p.6}\\
\K{Error} & 9.33 & \Sd{p.6}\\
\K{Error} & 9.15 & \Sd{p.6}\\
\K{Error} & 24.27 & \Sd{p.6}\\
\K{Error} & 7.38 & \Sd{p.6}\\
\K{Error} & 10.02 & \Sd{p.6}\\
\K{Error} & 7.76 & \Sd{p.6}\\
\K{Error} & 24.52 & \Sd{p.6}\\
\K{Error} & 7.46 & \Sd{p.6}\\
\K{Error} & 24.19 & \Sd{p.6}\\
\K{Error} & 7.40 & \Sd{p.6}\\
\K{Error} & 22.85 & \Sd{p.6}\\
\K{Error} & 6.71 & \Sd{p.6}\\
\K{Error} & 21.75 & \Sd{p.6}\\
\K{Error} & 6.05 & \Sd{p.6}\\
\K{Error} & 21.43 & \Sd{p.6}\\
\K{Error} & 5.71 & \Sd{p.6}\\
\K{Error} & 8.43 & \Sd{p.6}\\
\K{Error} & 7.89 & \Sd{p.6}\\
\K{Error} & 24.4 & \Sd{p.6}\\
\K{Error} & 7.1 & \Sd{p.6}\\
\K{Error} & 21.59 & \Sd{p.6}\\
\K{Error} & 21.99 & \Sd{p.6}\\
\K{Error} & 5.81 & \Sd{p.6}\\
\K{Error} & 21.84 & \Sd{p.6}\\
\K{Error} & 21.53 & \Sd{p.6}\\
\K{Error} & 5.60 & \Sd{p.6}\\
\K{Error} & 20.74 & \Sd{p.6}\\
\K{Error} & 5.25 & \Sd{p.6}\\
\K{Error} & 19.87 & \Sd{p.6}\\
\K{Error} & 4.60 & \Sd{p.6}\\
\K{Error} & 19.38 & \Sd{p.6}\\
\K{Error} & 4.49 & \Sd{p.6}\\
\K{Error} & 7.32 & \Sd{p.6}\\
\K{Error} & 6.66 & \Sd{p.6}\\
\K{Error} & 6.8 & \Sd{p.6}\\
\K{Error} & 4.94 & \Sd{p.6}\\
\K{Error} & 4.82 & \Sd{p.6}\\
\K{Error} & 3.57 & \Sd{p.6}\\
\K{Error} & 9.38 & \Sd{p.7}\\
\K{Error} & 8.81 & \Sd{p.7}\\
\K{Error} & 8.22 & \Sd{p.7}\\
\K{Error} & 8.39 & \Sd{p.7}\\
\K{Error} & 7.54 (7.72$\pm$0.16) & \Sd{p.7}\\
\K{Error} & 8.80 & \Sd{p.7}\\
\K{Error} & 8.75 & \Sd{p.7}\\
\K{Error} & 7.51 & \Sd{p.7}\\
\K{Error} & 7.17 & \Sd{p.7}\\
\K{Error} & 6.97 & \Sd{p.7}\\
\K{Error} & 6.43 (6.61$\pm$0.16) & \Sd{p.7}\\
\K{Error} & 7.93 & \Sd{p.7}\\
\K{Error} & 4.49\% & \Sd{p.7}\\
\K{Error} & 6.43\% & \Sd{p.8}\\
\K{Error} & 7.93\% & \Sd{p.8}\\
\K{Error} & 60\% & \Sd{p.8}\\
\K{Error} & 0.1\% & \Sd{p.8}\\
\K{Error} & 33.1\% & \Sd{p.12}\\
\K{Error} & 13.3\% & \Sd{p.12}\\
\K{Error} & 11.7\% & \Sd{p.12}\\
\K{Error} & 14.4\% & \Sd{p.12}\\
\K{Error} & 10.6\% & \Sd{p.12}\\
\K{Error} & 9.0\% & \Sd{p.12}\\
\K{Error} & 4.6\% & \Sd{p.12}\\
\K{Error} & 33.1 & \Sd{p.12}\\
\K{Error} & 13.3 & \Sd{p.12}\\
\K{Error} & 11.7 & \Sd{p.12}\\
\K{Error} & 14.4 & \Sd{p.12}\\
\K{Error} & 10.6 & \Sd{p.12}\\
\K{Error} & 8.9 & \Sd{p.12}\\
\K{Error} & 30.0 & \Sd{p.12}\\
\K{Error} & 29.9 & \Sd{p.12}\\
\K{Error} & 26.7 & \Sd{p.12}\\
\K{Error} & 26.9 & \Sd{p.12}\\
\K{Error} & 25.3 & \Sd{p.12}\\
\K{Error} & 9.0 & \Sd{p.12}\\
\K{Error} & 80\% & \Sd{p.7}\\
\K{Error} & 90\% & \Sd{p.7}\\
\end{kvtable}

\begin{kvtable}{mAP}{tab:sec_results-map}
\K{mAP} & 73.2 & \Sd{p.8}\\
\K{mAP} & 70.4 & \Sd{p.8}\\
\K{mAP} & 76.4 & \Sd{p.8}\\
\K{mAP} & 73.8 & \Sd{p.8}\\
\K{mAP} & 41.5 & \Sd{p.8}\\
\K{mAP} & 21.2 & \Sd{p.8}\\
\K{mAP} & 48.4 & \Sd{p.8}\\
\K{mAP} & 27.2 & \Sd{p.8}\\
\K{mAP} & 49.9 & \Sd{p.11}\\
\K{mAP} & 29.9 & \Sd{p.11}\\
\K{mAP} & 51.1 & \Sd{p.11}\\
\K{mAP} & 30.0 & \Sd{p.11}\\
\K{mAP} & 53.3 & \Sd{p.11}\\
\K{mAP} & 32.2 & \Sd{p.11}\\
\K{mAP} & 53.8 & \Sd{p.11}\\
\K{mAP} & 32.5 & \Sd{p.11}\\
\K{mAP} & 55.7 & \Sd{p.11}\\
\K{mAP} & 34.9 & \Sd{p.11}\\
\K{mAP} & 59.0 & \Sd{p.11}\\
\K{mAP} & 37.4 & \Sd{p.11}\\
\K{mAP} & 43.9 & \Sd{p.11}\\
\K{mAP} & 60.5 & \Sd{p.11}\\
\K{mAP} & 58.8 & \Sd{p.11}\\
\K{mAP} & 63.6 & \Sd{p.11}\\
\K{mAP} & 62.1 & \Sd{p.11}\\
\K{mAP} & 85.6 & \Sd{p.11}\\
\K{mAP} & 83.8 & \Sd{p.11}\\
\end{kvtable}

\begin{kvtable}{AP}{tab:sec_results-ap}
\K{AP} & 76.5 & \Sd{p.11}\\
\K{AP} & 79.0 & \Sd{p.11}\\
\K{AP} & 70.9 & \Sd{p.11}\\
\K{AP} & 65.5 & \Sd{p.11}\\
\K{AP} & 52.1 & \Sd{p.11}\\
\K{AP} & 83.1 & \Sd{p.11}\\
\K{AP} & 84.7 & \Sd{p.11}\\
\K{AP} & 86.4 & \Sd{p.11}\\
\K{AP} & 52.0 & \Sd{p.11}\\
\K{AP} & 81.9 & \Sd{p.11}\\
\K{AP} & 65.7 & \Sd{p.11}\\
\K{AP} & 84.8 & \Sd{p.11}\\
\K{AP} & 84.6 & \Sd{p.11}\\
\K{AP} & 77.5 & \Sd{p.11}\\
\K{AP} & 76.7 & \Sd{p.11}\\
\K{AP} & 38.8 & \Sd{p.11}\\
\K{AP} & 73.6 & \Sd{p.11}\\
\K{AP} & 73.9 & \Sd{p.11}\\
\K{AP} & 83.0 & \Sd{p.11}\\
\K{AP} & 72.6 & \Sd{p.11}\\
\K{AP} & 79.8 & \Sd{p.11}\\
\K{AP} & 80.7 & \Sd{p.11}\\
\K{AP} & 76.2 & \Sd{p.11}\\
\K{AP} & 68.3 & \Sd{p.11}\\
\K{AP} & 55.9 & \Sd{p.11}\\
\K{AP} & 85.1 & \Sd{p.11}\\
\K{AP} & 85.3 & \Sd{p.11}\\
\K{AP} & 89.8 & \Sd{p.11}\\
\K{AP} & 56.7 & \Sd{p.11}\\
\K{AP} & 87.8 & \Sd{p.11}\\
\K{AP} & 69.4 & \Sd{p.11}\\
\K{AP} & 88.3 & \Sd{p.11}\\
\K{AP} & 88.9 & \Sd{p.11}\\
\K{AP} & 80.9 & \Sd{p.11}\\
\K{AP} & 78.4 & \Sd{p.11}\\
\K{AP} & 41.7 & \Sd{p.11}\\
\K{AP} & 78.6 & \Sd{p.11}\\
\K{AP} & 72.0 & \Sd{p.11}\\
\K{AP} & 90.0 & \Sd{p.11}\\
\K{AP} & 89.6 & \Sd{p.11}\\
\K{AP} & 80.8 & \Sd{p.11}\\
\K{AP} & 76.1 & \Sd{p.11}\\
\K{AP} & 89.9 & \Sd{p.11}\\
\K{AP} & 75.5 & \Sd{p.11}\\
\K{AP} & 90.3 & \Sd{p.11}\\
\K{AP} & 89.1 & \Sd{p.11}\\
\K{AP} & 88.7 & \Sd{p.11}\\
\K{AP} & 65.4 & \Sd{p.11}\\
\K{AP} & 88.1 & \Sd{p.11}\\
\K{AP} & 85.6 & \Sd{p.11}\\
\K{AP} & 89.0 & \Sd{p.11}\\
\K{AP} & 86.8 & \Sd{p.11}\\
\K{AP} & 84.9 & \Sd{p.11}\\
\K{AP} & 74.3 & \Sd{p.11}\\
\K{AP} & 53.9 & \Sd{p.11}\\
\K{AP} & 49.8 & \Sd{p.11}\\
\K{AP} & 75.9 & \Sd{p.11}\\
\K{AP} & 88.5 & \Sd{p.11}\\
\K{AP} & 45.6 & \Sd{p.11}\\
\K{AP} & 77.1 & \Sd{p.11}\\
\K{AP} & 55.3 & \Sd{p.11}\\
\K{AP} & 86.9 & \Sd{p.11}\\
\K{AP} & 81.7 & \Sd{p.11}\\
\K{AP} & 79.6 & \Sd{p.11}\\
\K{AP} & 40.1 & \Sd{p.11}\\
\K{AP} & 60.9 & \Sd{p.11}\\
\K{AP} & 81.2 & \Sd{p.11}\\
\K{AP} & 61.5 & \Sd{p.11}\\
\K{AP} & 86.5 & \Sd{p.11}\\
\K{AP} & 81.6 & \Sd{p.11}\\
\K{AP} & 77.2 & \Sd{p.11}\\
\K{AP} & 58.0 & \Sd{p.11}\\
\K{AP} & 51.0 & \Sd{p.11}\\
\K{AP} & 76.6 & \Sd{p.11}\\
\K{AP} & 93.2 & \Sd{p.11}\\
\K{AP} & 48.6 & \Sd{p.11}\\
\K{AP} & 80.4 & \Sd{p.11}\\
\K{AP} & 59.0 & \Sd{p.11}\\
\K{AP} & 92.1 & \Sd{p.11}\\
\K{AP} & 48.1 & \Sd{p.11}\\
\K{AP} & 77.3 & \Sd{p.11}\\
\K{AP} & 66.5 & \Sd{p.11}\\
\K{AP} & 65.6 & \Sd{p.11}\\
\K{AP} & 88.4 & \Sd{p.11}\\
\K{AP} & 71.4 & \Sd{p.11}\\
\K{AP} & 86.3 & \Sd{p.11}\\
\K{AP} & 94.2 & \Sd{p.11}\\
\K{AP} & 66.8 & \Sd{p.11}\\
\K{AP} & 89.4 & \Sd{p.11}\\
\K{AP} & 69.2 & \Sd{p.11}\\
\K{AP} & 93.9 & \Sd{p.11}\\
\K{AP} & 91.9 & \Sd{p.11}\\
\K{AP} & 90.9 & \Sd{p.11}\\
\K{AP} & 67.9 & \Sd{p.11}\\
\K{AP} & 88.2 & \Sd{p.11}\\
\K{AP} & 76.8 & \Sd{p.11}\\
\K{AP} & 80.0 & \Sd{p.11}\\
\end{kvtable}

\begin{kvtable}{Class}{tab:sec_results-class}
\K{Class} & aero & \Sd{p.11}\\
\K{Class} & bike & \Sd{p.11}\\
\K{Class} & bird & \Sd{p.11}\\
\K{Class} & boat & \Sd{p.11}\\
\K{Class} & bottle & \Sd{p.11}\\
\K{Class} & bus & \Sd{p.11}\\
\K{Class} & car & \Sd{p.11}\\
\K{Class} & cat & \Sd{p.11}\\
\K{Class} & chair & \Sd{p.11}\\
\K{Class} & cow & \Sd{p.11}\\
\K{Class} & table & \Sd{p.11}\\
\K{Class} & dog & \Sd{p.11}\\
\K{Class} & horse & \Sd{p.11}\\
\K{Class} & mbike & \Sd{p.11}\\
\K{Class} & person & \Sd{p.11}\\
\K{Class} & plant & \Sd{p.11}\\
\K{Class} & sheep & \Sd{p.11}\\
\K{Class} & sofa & \Sd{p.11}\\
\K{Class} & train & \Sd{p.11}\\
\K{Class} & tv & \Sd{p.11}\\
\end{kvtable}

\begin{kvtable}{Improvement}{tab:sec_results-improvement}
\K{Improvement} & 28\% & \Sd{p.1}\\
\K{Improvement} & 18\% & \Sd{p.3}\\
\K{Improvement} & 2.8\% & \Sd{p.5}\\
\K{Improvement} & 3.5\% & \Sd{p.6}\\
\K{Improvement} & 6.0\% & \Sd{p.8}\\
\K{Improvement} & 6\% & \Sd{p.10}\\
\K{Improvement} & 6.9\% & \Sd{p.10}\\
\K{Improvement} & >3\% & \Sd{p.10}\\
\K{Improvement} & about 2 points & \Sd{p.10}\\
\K{Improvement} & about 1 point & \Sd{p.10}\\
\K{Improvement} & over 2 points & \Sd{p.11}\\
\K{Improvement} & 10 points & \Sd{p.11}\\
\K{Improvement} & 8.5 points & \Sd{p.12}\\
\K{Improvement} & 64\% & \Sd{p.12}\\
\end{kvtable}

\begin{kvtable}{Comparison}{tab:sec_results-comparison}
\K{Comparison} & 8x & \Sd{p.1}\\
\end{kvtable}

\begin{kvtable}{Competition}{tab:sec_results-competition}
\K{Competition} & ILSVRC 2015 & \Sd{p.1}\\
\K{Competition} & ILSVRC \& COCO 2015 & \Sd{p.1}\\
\K{Competition} & ILSVRC'14 & \Sd{p.6}\\
\K{Competition} & ILSVRC'15 & \Sd{p.6}\\
\K{Competition} & COCO 2015 & \Sd{p.11}\\
\K{Competition} & ILSVRC'13 & \Sd{p.12}\\
\K{Competition} & ILSVRC 14 & \Sd{p.12}\\
\end{kvtable}

\begin{kvtable}{Placement}{tab:sec_results-placement}
\K{Placement} & 1st place & \Sd{p.1}\\
\K{Placement} & 1st places & \Sd{p.1}\\
\K{Placement} & second place & \Sd{p.12}\\
\end{kvtable}


\end{document}
